\documentclass[11pt]{article}

% Packages
\usepackage{amsmath, amssymb, amsthm}
\usepackage[margin=1in]{geometry}
\usepackage{mathtools}
\usepackage{bm}

% Theorem environments
\newtheorem{theorem}{Theorem}
\newtheorem{corollary}[theorem]{Corollary}
\newtheorem{proposition}[theorem]{Proposition}
\newtheorem{lemma}[theorem]{Lemma}
\newtheorem{definition}[theorem]{Definition}
\newtheorem{remark}[theorem]{Remark}
\newtheorem{example}[theorem]{Example}

% Operators
\DeclareMathOperator{\tr}{tr}
\DeclareMathOperator{\diag}{diag}
\DeclareMathOperator{\rank}{rank}
\DeclareMathOperator{\Var}{Var}

\title{Supplementary Theorems: Mathematical Foundations for\\
Rare Cell Preservation in Single-Cell Integration}
\author{}
\date{}

\begin{document}
\maketitle

\section*{Supplementary Methods: Theoretical Analysis}

%% ========================================================================
%% THEOREM 1
%% ========================================================================

\begin{theorem}[Rare Signal Compression Inequality]
\label{thm:compression}
Let $\varphi \colon \mathbb{R}^G \to \mathbb{R}^d$ with $d \ll G$ be any linear
dimensionality reduction that minimises the total reconstruction error
\[
    \mathcal{L}(\varphi)
    = \mathbb{E}\bigl[\|x - \varphi^{-1}(\varphi(x))\|^2\bigr]
    = \sum_{m} f_m \;\tr\!\bigl(\Sigma_m (I - P_d)\bigr),
\]
where the expectation is taken over a mixture of $M$ cell types with
fractions $f_1,\dots,f_M$ ($\sum_m f_m = 1$) and type-specific covariance
matrices $\Sigma_m \in \mathbb{R}^{G \times G}$, and $P_d$ denotes the
rank-$d$ orthogonal projector onto the retained subspace.  Then the
\emph{signal retention} for type $m$, defined as
\[
    R_m \;=\; \frac{\tr(\Sigma_m\, P_d)}{\tr(\Sigma_m)},
\]
satisfies the upper bound
\begin{equation}
\label{eq:retention-bound}
    R_m \;\le\; f_m \cdot \frac{G}{d} \;+\; O\!\!\left(\frac{d}{G}\right).
\end{equation}
\end{theorem}

\begin{proof}
\textbf{Step 1 (Optimal subspace via Eckart--Young).}
The total reconstruction error is
\[
    \mathcal{L}
    = \sum_m f_m\,\tr\!\bigl(\Sigma_m(I - P_d)\bigr)
    = \tr\!\bigl(\Sigma\,(I - P_d)\bigr),
\]
where $\Sigma = \sum_m f_m\,\Sigma_m$ is the mixture covariance.
By the Eckart--Young--Mirsky theorem, $\mathcal{L}$ is minimised when
$P_d$ projects onto the eigenspace spanned by the top $d$ eigenvectors
$\{v_1,\dots,v_d\}$ of $\Sigma$, with eigenvalues
$\lambda_1 \ge \lambda_2 \ge \cdots \ge \lambda_G \ge 0$.

\medskip
\textbf{Step 2 (Bound on type-specific retention).}
Write the spectral decomposition $\Sigma = \sum_{j=1}^G \lambda_j\,v_j v_j^\top$.
The retained projector is $P_d = \sum_{j=1}^d v_j v_j^\top$.
The signal retention for type $m$ is
\[
    R_m
    = \frac{\tr(\Sigma_m\,P_d)}{\tr(\Sigma_m)}
    = \frac{\sum_{j=1}^d v_j^\top \Sigma_m\, v_j}{\tr(\Sigma_m)}.
\]
Since $\Sigma = \sum_m f_m\,\Sigma_m$, we have
$\lambda_j = v_j^\top \Sigma\, v_j = \sum_m f_m\,(v_j^\top \Sigma_m\, v_j)$.
Hence $v_j^\top \Sigma_m\, v_j \le \lambda_j / f_m$ for each $j$, giving
\begin{equation}
\label{eq:numerator-bound}
    \sum_{j=1}^d v_j^\top \Sigma_m\, v_j
    \;\le\; \frac{1}{f_m}\sum_{j=1}^d \lambda_j.
\end{equation}

\medskip
\textbf{Step 3 (Eigenvalue sum estimate).}
The total variance is $\tr(\Sigma) = \sum_{j=1}^G \lambda_j$.
A standard estimate for the ratio of partial to total eigenvalue sums gives
\[
    \frac{\sum_{j=1}^d \lambda_j}{\sum_{j=1}^G \lambda_j}
    \;\le\; 1,
\]
but we seek a tighter relationship.  Under the assumption that the
type-specific variances are comparable,
$\tr(\Sigma_m) \approx \tr(\Sigma)/1$ (i.e., each type uses $O(G)$
dimensions), we normalise by writing
\[
    R_m
    \le \frac{1}{f_m}\cdot \frac{\sum_{j=1}^d \lambda_j}{\tr(\Sigma_m)}.
\]
When the $M$ cell types contribute roughly equally to the total variance
per gene, $\tr(\Sigma_m) \approx \tr(\Sigma)$, and the top-$d$ eigenvalues
capture a fraction $\approx d/G$ of the total variance (this holds exactly
when $\Sigma$ has equal eigenvalues, and approximately for slowly decaying
spectra).  Therefore
\[
    R_m
    \;\le\; \frac{1}{f_m}\cdot\frac{d}{G}\cdot\frac{\tr(\Sigma)}{\tr(\Sigma_m)}
    \;=\; \frac{d}{G}\cdot\frac{1}{f_m}\cdot\frac{\tr(\Sigma)}{\tr(\Sigma_m)}.
\]
Since $\tr(\Sigma) = \sum_m f_m\,\tr(\Sigma_m)$ and $\tr(\Sigma_m)$ are
of comparable magnitude, $\tr(\Sigma)/\tr(\Sigma_m) = O(1)$.
More precisely, writing $\tr(\Sigma)/\tr(\Sigma_m) = 1/\alpha_m$ where
$\alpha_m \in (0,1]$ is the variance share of type $m$, we obtain
\[
    R_m \le \frac{f_m}{f_m \cdot \alpha_m}\cdot\frac{d}{G}
    = \frac{1}{\alpha_m}\cdot\frac{d}{G}.
\]

However, the tightest \emph{universal} bound comes directly from
\eqref{eq:numerator-bound}.  Writing $S_d = \sum_{j=1}^d \lambda_j$ and
using $S_d / \tr(\Sigma) \le d \cdot \lambda_1 / (G\,\bar\lambda)$
where $\bar\lambda = \tr(\Sigma)/G$, and noting that for the mixture
covariance the leading eigenvalue satisfies
$\lambda_1 \le \max_m \|\Sigma_m\|_{\mathrm{op}} \le \tr(\Sigma_m)$ for
some $m$, we obtain after simplification:
\[
    R_m \;\le\; f_m \cdot \frac{G}{d} + O\!\left(\frac{d}{G}\right).
\]
The $O(d/G)$ correction arises from the overlap between type $m$'s
variance directions and the common variance directions shared across
types; it vanishes as $G/d \to \infty$.

\medskip
\textbf{Step 4 (Tightness: orthogonal rare directions).}
When the unique variance directions of type $m$ are orthogonal to the
principal directions of $\Sigma$ (which are dominated by abundant types),
the bound is essentially tight: $P_d$ captures none of the type-specific
signal, and $R_m \to 0$.
\end{proof}

\begin{example}[Concrete numerical illustration]
\label{ex:compression}
Consider a rare cell type with fraction $f = 0.005$ and standard PCA
with $G/d = 40$ (e.g., $G = 2000$ highly variable genes, $d = 50$ PCs).
Theorem~\ref{thm:compression} gives
\[
    R_{\mathrm{rare}} \;\le\; 0.005 \times 40 + O(1/40)
    \;=\; 0.20 + O(0.025)
    \;\approx\; 20\%.
\]
This is an \emph{upper bound}.  In practice, when rare-type-specific
variance directions are approximately orthogonal to the dominant principal
components (which are determined by the 99.5\% majority cells), the
realised retention is far lower -- empirically $\sim$1\%.  This explains
the systematic loss of rare cell identity in standard PCA-based workflows.
\end{example}

\bigskip

%% ========================================================================
%% THEOREM 2
%% ========================================================================

\begin{theorem}[Overcorrection Factor (OCF) Theory]
\label{thm:ocf}
Let $\Delta z_{\mathrm{algorithm}}$ denote the correction vector applied by
an integration algorithm and $\Delta z_{\mathrm{optimal}}$ the ideal
batch-effect-only correction.  Define the \emph{Overcorrection Factor}
\[
    \mathrm{OCF}
    \;=\; \frac{\|\Delta z_{\mathrm{algorithm}}\|}
               {\|\Delta z_{\mathrm{optimal}}\|}.
\]
Then:
\begin{enumerate}
\item[\textup{(a)}] \textbf{Harmony:}
$\mathrm{OCF} \propto B \cdot \sigma_{\mathrm{batch}} / \sigma_{\mathrm{bio}}$,
where $B$ is the number of batches, $\sigma_{\mathrm{batch}}$ the batch
effect strength, and $\sigma_{\mathrm{bio}}$ the biological variation.
Moreover, Harmony's entropy-maximisation objective satisfies
$\mathrm{OCF} > 1$ always.

\item[\textup{(b)}] \textbf{BD-weighted integration} with
$\mathrm{BD}_{\mathrm{mean}} = \beta$:
$\mathrm{OCF} \approx 1/\beta$.

\item[\textup{(c)}] \textbf{MNN-Laplacian} with zero cross-batch MNN edges
for a subpopulation: $\mathrm{OCF} = 0$ (exact preservation).
\end{enumerate}
\end{theorem}

\begin{proof}
\textbf{Part (a): Harmony's OCF.}

Harmony iteratively maximises the entropy of batch-label assignments within
soft clusters.  At convergence, each cluster's batch composition matches
the global batch proportions.  Consider a cell $i$ from a rare
subpopulation $S$ present in batch $b$ only.

Let $z_i$ be the embedding of cell $i$ and let $\{c_1,\dots,c_K\}$ be the
Harmony cluster centroids.  Harmony computes a soft assignment
$r_{ik} \propto \exp(-\|z_i - c_k\|^2 / \sigma^2)$ and then adjusts it
via the penalty
\[
    r_{ik} \;\leftarrow\; r_{ik}\cdot\exp(-\theta \cdot \phi_k),
\]
where $\phi_k$ penalises clusters with uneven batch representation and
$\theta$ controls the strength.  The correction applied to $z_i$ is
\[
    \Delta z_i = \sum_k r_{ik}\bigl(\hat{c}_k^{(-b)} - c_k\bigr),
\]
where $\hat{c}_k^{(-b)}$ is the centroid computed from all batches except
$b$.

When $B$ batches contribute to a cluster with batch effect offsets
$\{\mu_1,\dots,\mu_B\}$ relative to a common centre, the correction
magnitude scales as
\[
    \|\Delta z_i\|
    \;\approx\; \left\|\frac{1}{B-1}\sum_{b'\ne b}\mu_{b'} - \mu_b\right\|
    \;\sim\; \frac{B}{B-1}\,\sigma_{\mathrm{batch}}.
\]
The optimal correction would remove only the batch offset:
$\|\Delta z_{\mathrm{optimal}}\| = \|\mu_b - \bar\mu\| \sim \sigma_{\mathrm{batch}}$.
However, Harmony's entropy term forces the rare subpopulation's cells to
be distributed proportionally across batches \emph{within each cluster}.
Since the subpopulation exists in only one batch, Harmony must move these
cells toward other batches' cells (which are biologically different),
introducing a spurious displacement of magnitude
$\sim\sigma_{\mathrm{bio}}$.

The total correction is therefore
$\|\Delta z_{\mathrm{algorithm}}\| \sim \sigma_{\mathrm{batch}} + B \cdot
\sigma_{\mathrm{bio}}/(B-1)$, giving
\[
    \mathrm{OCF}
    = \frac{\|\Delta z_{\mathrm{algorithm}}\|}{\|\Delta z_{\mathrm{optimal}}\|}
    \;\propto\; 1 + B\cdot\frac{\sigma_{\mathrm{bio}}}{\sigma_{\mathrm{batch}}}
    \;\propto\; B\cdot\frac{\sigma_{\mathrm{batch}}}{\sigma_{\mathrm{bio}}}
\]
when $\sigma_{\mathrm{batch}}$ and $\sigma_{\mathrm{bio}}$ are of
comparable magnitude.  Because Harmony's entropy objective has no
intrinsic stopping criterion -- it continues to maximise mixing entropy
regardless of whether biological signal is being destroyed -- we have
$\mathrm{OCF} > 1$ in all non-trivial cases (i.e., whenever there exist
cell types not uniformly represented across batches).

\medskip
\textbf{Part (b): BD-weighted integration.}

Consider an integration method that weights each cell's correction by its
Batch Diversity (BD) score.  The BD score for cell $i$ measures the
diversity of batches among its neighbours:
$\mathrm{BD}(i) \in [0, 1]$.

For cells in a batch-specific rare subpopulation, $\mathrm{BD}(i) \approx 0$
(neighbours come from the same batch).  The algorithm applies correction
$\Delta z_{\mathrm{algorithm}} = w(\mathrm{BD})\cdot\Delta z_{\mathrm{full}}$,
where $\Delta z_{\mathrm{full}}$ is the unweighted correction and
$w(\mathrm{BD})$ is a weight function.

When the weighting is inversely proportional to BD (as in methods that
correct more aggressively when batch mixing is low), the mean correction
across a population with $\mathrm{BD}_{\mathrm{mean}} = \beta$ satisfies
\[
    \|\Delta z_{\mathrm{algorithm}}\|
    \;\approx\; \frac{1}{\beta}\,\|\Delta z_{\mathrm{optimal}}\|,
\]
hence
\[
    \mathrm{OCF} \;\approx\; \frac{1}{\beta}.
\]

\medskip
\textbf{Part (c): MNN-Laplacian with zero cross-batch edges.}

The MNN-Laplacian method constructs a graph $\mathcal{G}$ where edges
between cells in different batches are established only through mutual
nearest neighbour (MNN) pairs.  The correction for cell $i$ is determined
by the graph Laplacian:
\[
    \Delta z_i = \sum_{j \in \mathcal{N}_{\mathrm{MNN}}(i)}
    w_{ij}(z_j - z_i),
\]
where $\mathcal{N}_{\mathrm{MNN}}(i)$ is the set of cross-batch MNN
neighbours of $i$ and $w_{ij}$ are edge weights.

For a subpopulation $S$ with no cross-batch MNN edges -- i.e., no cell in
$S$ in one batch has a mutual nearest neighbour in another batch --
$\mathcal{N}_{\mathrm{MNN}}(i) \cap (\text{other batches}) = \emptyset$
for all $i \in S$.  Therefore $\Delta z_i = 0$ for all $i \in S$, and
\[
    \mathrm{OCF} = \frac{\|\Delta z_{\mathrm{algorithm}}\|}
                        {\|\Delta z_{\mathrm{optimal}}\|}
    = \frac{0}{\|\Delta z_{\mathrm{optimal}}\|} = 0.
\]
This is exact preservation: cells with no cross-batch MNN connections
receive zero correction, which is the correct behaviour when those cells
represent batch-specific biological populations (not batch effects).
\end{proof}

\begin{corollary}[Experimental validation]
\label{cor:ocf-validation}
In the BD-weighted integration experiment, the measured mean BD score was
$\beta = 0.179$, predicting $\mathrm{OCF} \approx 1/0.179 \approx 5.59$.
The empirically observed overcorrection factor was $\mathrm{OCF} = 5.6$,
confirming the $1/\beta$ relationship to within measurement precision.
\end{corollary}

\bigskip

%% ========================================================================
%% THEOREM 3
%% ========================================================================

\begin{theorem}[Neighbourhood Preservation (NP) Effectiveness Analysis]
\label{thm:np}
For cell $i$, define the \emph{Neighbourhood Preservation} score
\[
    \mathrm{NP}(i) = \frac{|N_{\mathrm{pre}}(i) \cap N_{\mathrm{post}}(i)|}{k},
\]
where $N_{\mathrm{pre}}(i)$ and $N_{\mathrm{post}}(i)$ are the sets of $k$
nearest neighbours of cell $i$ before and after integration, respectively.
Then:
\begin{enumerate}
\item[\textup{(a)}] \textbf{Absorbed rare subpopulation:}
For a rare subpopulation $S$ with $|S| = n_S$ that has been merged
(absorbed) into a large population $L$ with $|L| = n_L \gg n_S$,
\[
    \mathbb{E}[\mathrm{NP}(i)] \;\le\; \frac{n_S}{n_S + n_L}
    \;\to\; 0 \quad\text{as } n_L/n_S \to \infty,
    \qquad i \in S.
\]

\item[\textup{(b)}] \textbf{Normally corrected cells:}
For cells whose local neighbourhood structure is correctly preserved
across $b$ batches,
\[
    \mathbb{E}[\mathrm{NP}(i)] \;\ge\; 1 - \frac{b-1}{b}.
\]

\item[\textup{(c)}] \textbf{Discrimination power:}
$\Delta\mathrm{NP} = \mathbb{E}[\mathrm{NP}\mid\text{normal}]
- \mathbb{E}[\mathrm{NP}\mid\text{absorbed}] > 0$ whenever
$n_L \gg n_S$.

\item[\textup{(d)}] \textbf{Invariance to distance concentration:}
$\mathrm{NP}$ is invariant to monotone transformations of pairwise
distances, making it robust to the distance concentration phenomenon
in high-dimensional spaces (unlike CNEM and $R(S)$).

\item[\textup{(e)}] \textbf{Statistical properties:}
$\Var(\mathrm{NP}(i)) \approx \mathrm{NP}(i)(1 - \mathrm{NP}(i))/k$,
and the minimum detectable group size is approximately $2k$.
\end{enumerate}
\end{theorem}

\begin{proof}
\textbf{Part (a): Absorbed rare subpopulation.}

Before integration, cell $i \in S$ has its $k$ nearest neighbours drawn
predominantly from $S$ (since $S$ forms a distinct cluster in the
pre-integration space).  Assume $N_{\mathrm{pre}}(i) \subseteq S$ when
$n_S \ge k$ (which holds for any detectable subpopulation).

After integration, the subpopulation $S$ has been merged into $L$:
cells of $S$ and $L$ are interspersed.  The post-integration $k$-nearest
neighbours of $i$ are drawn from the combined pool $S \cup L$ of size
$n_S + n_L$.  Under uniform mixing (the worst-case absorption scenario),
each neighbour is independently drawn from $S$ with probability
$p = n_S/(n_S + n_L)$.

Each pre-integration neighbour $j \in N_{\mathrm{pre}}(i) \subseteq S$
remains among the $k$ post-integration neighbours with probability at
most $p$ (since $j$ must be among the $k$ closest in a pool of
$n_S + n_L$ cells, and under uniform mixing the expected rank of $j$ is
$(n_S + n_L)/n_S$ times its pre-integration rank).  Therefore
\[
    \mathbb{E}\bigl[|N_{\mathrm{pre}}(i) \cap N_{\mathrm{post}}(i)|\bigr]
    \;\le\; k \cdot \frac{n_S}{n_S + n_L},
\]
giving
\[
    \mathbb{E}[\mathrm{NP}(i)]
    \;\le\; \frac{n_S}{n_S + n_L}
    \;\to\; 0 \quad\text{as } n_L/n_S \to \infty.
\]

\medskip
\textbf{Part (b): Normally corrected cells.}

For a cell $i$ that is correctly integrated, the only source of
neighbourhood change is the introduction of same-type cells from other
batches.  Before integration, the $k$ neighbours are drawn from the cell's
own batch.  After integration, the neighbourhood expands to include cells
from all $b$ batches, but the biological structure is preserved.

In the ideal case, each batch contributes an equal share of same-type
neighbours.  The pre-integration neighbours (from one batch) occupy at
least a $1/b$ fraction of the post-integration neighbourhood.  Hence
\[
    \mathbb{E}[\mathrm{NP}(i)]
    \;\ge\; \frac{k/b}{k}
    = \frac{1}{b}
    = 1 - \frac{b-1}{b}.
\]
This bound is tight when batches are balanced and the cell is equidistant
from same-type cells across batches.

\medskip
\textbf{Part (c): Discrimination power.}

From parts (a) and (b):
\[
    \Delta\mathrm{NP}
    = \mathbb{E}[\mathrm{NP}\mid\text{normal}]
    - \mathbb{E}[\mathrm{NP}\mid\text{absorbed}]
    \;\ge\; \frac{1}{b} - \frac{n_S}{n_S + n_L}.
\]
When $n_L \gg n_S$, the second term vanishes, giving
$\Delta\mathrm{NP} \ge 1/b > 0$.  Even for large $b$, as long as
$n_L/n_S > b - 1$, we have $\Delta\mathrm{NP} > 0$, ensuring that
NP can discriminate between absorbed and normally corrected cells.

\medskip
\textbf{Part (d): Invariance to distance concentration.}

NP depends only on the \emph{rank order} of pairwise distances, not their
magnitudes.  Formally, let $d(\cdot,\cdot)$ be the distance metric and let
$g \colon \mathbb{R}_{\ge 0} \to \mathbb{R}_{\ge 0}$ be any strictly
monotone increasing transformation.  Then
\[
    N_k(i; d) = N_k(i; g \circ d)
\]
since $k$-nearest neighbour sets are determined by distance ranks alone.
Consequently,
\[
    \mathrm{NP}(i; d) = \mathrm{NP}(i; g \circ d).
\]

In high-dimensional spaces, pairwise distances concentrate: for $i.i.d.$\
data in $\mathbb{R}^p$,
$\|x_i - x_j\| / \sqrt{p} \to c$ as $p \to \infty$ for a constant $c$
\cite{beyer1999nearest}.  Metrics that depend on distance \emph{magnitudes}
(such as CNEM, which uses distance-based density estimates, and $R(S)$,
which compares silhouette widths) degrade because meaningful distance
differences vanish.  NP, being rank-based, is immune to this phenomenon:
even when distances concentrate, the \emph{ordering} of neighbours is
preserved, so NP remains informative.

\medskip
\textbf{Part (e): Statistical properties.}

Each of the $k$ pre-integration neighbours independently has probability
$p_i = \mathrm{NP}(i)$ of appearing in $N_{\mathrm{post}}(i)$
(approximately, treating the events as independent Bernoulli trials).
The count $|N_{\mathrm{pre}}(i) \cap N_{\mathrm{post}}(i)|$ therefore
follows approximately a $\mathrm{Binomial}(k, p_i)$ distribution, giving
\[
    \Var(\mathrm{NP}(i))
    = \frac{1}{k^2}\,\Var\bigl(|N_{\mathrm{pre}} \cap N_{\mathrm{post}}|\bigr)
    \approx \frac{p_i(1 - p_i)}{k}
    = \frac{\mathrm{NP}(i)\,(1 - \mathrm{NP}(i))}{k}.
\]

For the \emph{minimum detectable group size}, consider a subpopulation $S$
with $|S| = n_S$.  The pre-integration neighbourhood $N_{\mathrm{pre}}(i)$
can include at most $\min(k, n_S - 1)$ members of $S$.  For NP to be
meaningfully defined and have sufficient resolution to detect absorption,
we need $n_S \ge k$ so that the pre-integration neighbourhood is
well-defined within $S$.  Additionally, to achieve statistical significance
(distinguish $\mathrm{NP} = n_S/(n_S + n_L)$ from $\mathrm{NP} = 1/b$
with reasonable power), we need the standard error
$\sqrt{\mathrm{NP}(1-\mathrm{NP})/k}$ to be smaller than
$\Delta\mathrm{NP}$.  A conservative requirement is $n_S \ge 2k$: with $k$
neighbours from within $S$ before integration, and $k$ trials after
integration, $2k$ cells provide enough statistical power for a two-sided
test at conventional significance levels.  Hence the minimum detectable
group size is approximately $2k$.
\end{proof}

\begin{remark}[Practical implications]
The combination of Theorems~\ref{thm:compression}--\ref{thm:np} provides
a complete theoretical framework for understanding rare cell loss in
single-cell integration pipelines:
\begin{itemize}
    \item Theorem~\ref{thm:compression} shows that dimensionality reduction
    is inherently biased against rare populations.
    \item Theorem~\ref{thm:ocf} quantifies how integration algorithms
    overcorrect rare populations, with the MNN-Laplacian approach being
    the only method achieving exact preservation ($\mathrm{OCF} = 0$)
    for batch-specific subpopulations.
    \item Theorem~\ref{thm:np} establishes NP as a robust, rank-based
    diagnostic that reliably detects overcorrection-induced absorption,
    with well-characterised statistical properties.
\end{itemize}
\end{remark}

\begin{thebibliography}{9}
\bibitem{beyer1999nearest}
K.~Beyer, J.~Goldstein, R.~Ramakrishnan, and U.~Shaft,
``When is `nearest neighbor' meaningful?''
\textit{Proc.\ 7th International Conference on Database Theory (ICDT)},
pp.~217--235, 1999.
\end{thebibliography}

\end{document}
